% Ad Familiares 13.66 (ad-familiares-13-66)
% Source: https://raw.githubusercontent.com/PerseusDL/canonical-latinLit/master/data/phi0474/phi056/phi0474.phi056.perseus-lat1.xml
% Fetched by scripts/fetch_latin.py

% Dateline (Perseus): Scr. Romae vel ex. m. Ian. vel in. Febr. a. 709 (45) .

\ciceroLetterOpener{M. CICERO P. SERVILIO S.}

\ciceroSection{1} A. Caecinam, maxime proprium clientem familiae vestrae, non commendarem tibi, cum scirem qua fide in tuos, qua clementia in calamitosos soleres esse, nisi me et patris eius, quo sum familiarissime usus, memoria et huius fortuna ita moveret, ut hominis omnibus mecum studiis officiisque coniunctissimi movere debebat. A te hoc omni contentione peto, sic ut maiore cura, maiore animi labore petere non possim, ut ad ea, quae tua sponte sine cuiusquam commendatione faceres in hominem tantum et talem et calamitosum, aliquem adferant cumulum meae litterae, quo studiosius eum quibuscumque rebus possis iuves.

\ciceroSection{2} quod si Romae fuisses, etiam salutem A. Caecinae essemus, ut opinio mea fert, per te consecuti ; de qua tamen magnam spem habemus freti clementia conlegae tui. nunc quoniam tuam iustitiam secutus tutissimum sibi portum provinciam istam duxit esse, etiam atque etiam te rogo atque oro ut eum et in reliquus veteris negotiationis conligendis iuves et ceteris rebus tegas atque tueare. hoc mihi gratius facere nihil potes.
